\section{从极端客户到通用平台：Inductive Productization}

企业软件常遇到一种反直觉情况：只有一个客户提出某个问题，并不代表问题不重要；这个客户可能比市场提前五年进入某个约束。讲者把 enterprise demand 描述成 power law，而不是所有需求围绕平均值分布。Forward-deployed 团队先在真实 workflow 中解决极端案例，再通过重复案例识别 invariant，最终形成产品接口。

\subsection{不是定制开发，也不是脱离现场的产品规划}

纯定制开发会把每个客户差异写进分支，无法积累；脱离现场的产品规划则容易把最有价值的极端约束平均掉。Inductive productization 介于两者之间：第一次记录完整情境，第二次寻找结构相似性，第三次才开始形成稳定 abstraction。Operation Warp Speed 的快速响应被讲者归因于更早在油气生产优化中解决过结构相似的资源协调问题。

\begin{figure}[H]
\centering
\includegraphics[width=0.94\textwidth]{images/07-inductive-productization.png}
\caption{从单个未来客户到通用产品：重复、分离局部细节、编码 invariant。}
\figsource{依据字幕 19:00--21:30 重绘，`images/07-inductive-productization.png`。}
\end{figure}

\begin{knowledgebox}{读图：什么才算 ``structurally identical''}
两个行业可以使用完全不同名词，却拥有相同约束图。例如都有稀缺资源、需求优先级、动态供应、审核权限和结果反馈，就可能共享 optimization 与 workflow primitive。相似不能只靠故事类比，需要把 entity、constraint、decision 和 outcome 明确映射。
\end{knowledgebox}

\begin{warningbox}{Forward-deployed 经验也会过拟合}
现场团队容易把最响亮客户的流程当成市场事实。避免过拟合需要记录哪些逻辑属于法规、组织习惯或临时 workaround，哪些是真正跨客户 invariant。产品化后还要验证新客户能否在不复制原团队人力的情况下采用。
\end{warningbox}

\subsection{本章小结}

极端任务环境可以提前暴露未来基础设施问题，但只有通过重复、对照和抽象才能转化为平台。Apollo/Rubix 本身就是这种路径的结果：先解决内部 fleet 的交付与运行，再逐步外化为通用能力。下一章把同一思路应用到 Project Maven，从视觉检测扩展到完整决策链。

